\subsubsection{Alcance del informe de estado de expedientes}\label{sec:diseno-informes}

La concreción de la \cref{tab:rf-19} en esta entrega aborda el estado de
expedientes y su carga de asignaciones; no agota el catálogo de tipos ni la
medición de desempeño. La \cref{tab:cu-16} conserva como alcance el aviso por
correo y la estimación de tiempo en su flujo alterno. La implementación ofrece
avisos internos y estados de procesamiento; el canal de correo y una estimación
validada de terminación permanecen pendientes. Estos límites no sustituyen
los requisitos ni permiten considerar el caso de uso completo.

El informe de estado representa una observación administrativa de los expedientes
originalmente creados dentro de un intervalo UTC. El inicio es inclusivo, el
fin exclusivo y la duración máxima es de 366 días. El estado activo o cerrado
se consulta al capturar los datos; no se reconstruye el estado que cada
expediente tenía al final del periodo. La carga de trabajo procede de las
asignaciones de esos mismos expedientes. Por ello, el resultado no es un censo
histórico, una medición de actividad durante el intervalo ni una calificación
de desempeño o efectividad jurídica.

Owner solicita el ámbito del despacho; Litigante, sus expedientes asignados.
Cada informe es privado de su solicitante, incluso frente a otro Owner.
Paralegal y Cliente no acceden a este agregado. La selección opcional de un
litigante utiliza el mismo criterio de pertenencia que el filtro: Owner ve
los litigantes activos, aunque carezcan de expedientes, y Litigante ve a los
activos con quienes comparte algún expediente, incluidos los cerrados. El
selector se pagina y no utiliza la carga de expedientes activos del tablero
como sustituto del directorio autorizado.

La solicitud se identifica con una operación repetible y una captura inmutable,
compartida por PDF y CSV. El trabajo puede continuar fuera de la petición HTTP,
pero su publicación y descarga vuelven a comprobar el acceso a todos los
expedientes capturados. La revocación de uno impide entregar el conjunto; no se
eliminan filas para presentar como completo un resultado distinto. Los avisos
de disponibilidad o fallo son internos y su lectura requiere acuse explícito.
La especificación concreta está en \rutarepo{docs/case-reports-api.md}; la
\cref{sec:impl-informes} describe su composición, y la
\cref{sec:pruebas-informes} separa implementación y aceptación reproducida.
